\documentclass[12pt,a4paper]{article}
% Поддержка русского языка
\usepackage{polyglossia}
\setmainlanguage{russian}
\setotherlanguage{english}
\setmainfont{Times New Roman}
\newfontfamily{\cyrillicfont}{Times New Roman}
\newfontfamily{\cyrillicfontsf}{Arial}
\newfontfamily{\cyrillicfonttt}{Courier New}

\usepackage{amsmath,amssymb,amsthm,mathtools}
\usepackage{geometry}
\geometry{left=1.5cm,right=1.5cm,top=1.5cm,bottom=2.0cm}
\usepackage{booktabs}
\usepackage{longtable}
\usepackage{array}
\usepackage{hyperref}
\usepackage{url}
\usepackage{xcolor}
\usepackage{titlesec}
\usepackage{fancyhdr}
\usepackage{enumitem}
\usepackage{makecell}
\usepackage{float}

\titleformat{\section}{\Large\bfseries}{\thesection}{1em}{}
\titleformat{\subsection}{\large\bfseries}{\thesubsection}{1em}{}

\hypersetup{colorlinks=true,linkcolor=blue,citecolor=blue,urlcolor=blue}

% --- Фундаментальные константы YUCT (оставляем как есть) ---
\newcommand{\REL}{\mathcal{K}}          % относительная эффективность координации
\newcommand{\ABS}{K_{\mathrm{eff}}}     % абсолютная
\newcommand{\kappac}{\kappa_c}
\newcommand{\betaf}{\beta}
\newcommand{\Sodd}{S_{\mathrm{odd}}}
\newcommand{\Seven}{S_{\mathrm{even}}}

\begin{document}
	
	\begin{titlepage}
		\centering
		\vspace*{2cm}
		{\Huge\bfseries Приложение BCLM-LL-MPS\\[0.3cm]
			Протокол валидации на нескольких видах\\для дизайна долгой жизни\par}
		\vspace{0.5cm}
		{\Large\bfseries Параллельный эксперимент на пяти видах млекопитающих\\для подтверждения координационной теории старения YUCT\par}
		\vspace{1.5cm}
		{\large A.\,V.\,Якушев\par}
		{\normalsize Исследования YUCT, \url{https://yuct.org/}\par}
		\vspace{0.5cm}
		{\normalsize Главный DOI: \url{https://doi.org/10.5281/zenodo.18444598}\par}
		\vspace{0.5cm}
		{\normalsize Июнь 2026 (v1.2 – межвидовая согласованность)\par}
		\vfill
		\begin{abstract}
			\noindent
			Биологический лагранжиан координации YUCT (BCLM) позволяет предсказывать эффективность координации \(K_{\mathrm{eff}}\) для любого белка или пути, исходя только из данных последовательности. Чтобы продемонстрировать универсальность координационной теории старения, мы представляем параллельный экспериментальный дизайн для пяти широко используемых лабораторных млекопитающих: мыши (\textit{Mus musculus}), крысы (\textit{Rattus norvegicus}), свиньи (\textit{Sus scrofa}), яванского макака (\textit{Macaca fascicularis}) и собаки (\textit{Canis familiaris}). Для каждого вида мы идентифицируем ортологи целевых генов, вычисляем их нативную и оптимизированную относительную эффективность координации \(\REL\) с помощью конвейера BCLM и выводим количественные предсказания для частоты мутаций, окислительного стресса, агрегации белков, скорости миграции и репликативной продолжительности жизни непосредственно из универсального закона ошибок \(\varepsilon = \kappa_c \alpha \REL^{-\beta}\) (\(\beta=2/3\), \(\kappa_c=1/3\)). Сравнительная таблица включает соответствующие значения для клеток человека (из сопутствующего протокола BCLM‑LL). Эксперимент опирается исключительно на коммерчески доступные клеточные линии, общедоступные базы данных последовательностей и стандартные методы клеточной биологии. Модуль контролируемой регрессии (CCR) подтверждает обратимость стареющего фенотипа у всех видов. Все предсказания предварительно зарегистрированы и фальсифицируемы. Протокол намеренно опускает методы \textit{in vivo} доставки; этические аспекты обсуждаются в Приложении~\(\Sigma\).
		\end{abstract}
	\end{titlepage}
	
	\tableofcontents
	\newpage
	
	\section{Введение}
	\label{sec:intro}
	Универсальный закон масштабирования ошибок YUCT,
	\begin{equation}
		\varepsilon = \kappa_c \, \alpha \, \REL^{-\beta}, \qquad \beta = \frac{2}{3},\quad \kappa_c = \frac{1}{3},
		\label{eq:error-law}
	\end{equation}
	применим к любой координированной системе, независимо от масштаба и субстрата. В клеточном контексте относительная эффективность \(\REL\) количественно определяет точность молекулярных машин; её возрастное снижение ускоряет накопление повреждений и в конечном счёте запускает сенесценс.
	
	Биологический лагранжиан координации (BCLM, Приложение BCLM) предоставляет единый метод для вычисления \(\REL\) по данным последовательности и для создания кодон-оптимизированных конструкций, которые повышают \(\REL\) до ювенильного уровня. Сопутствующий протокол BCLM‑LL описывает полный эксперимент для кардиомиоцитов человека. Здесь мы распространяем этот дизайн на пять дополнительных видов млекопитающих — мышь, крысу, свинью, макака и собаку, — каждый из которых имеет полностью аннотированный геном, охарактеризованное использование кодонов и коммерчески доступные первичные клеточные линии. Цель — проверить универсальность механизма старения YUCT: если одна и та же количественная связь между \(\REL\) и клеточным здоровым долголетием выполняется у всех видов, координационная теория старения получит сильную поддержку.
	
	\subsection{Существующая экспериментальная поддержка}
	Молекулярные мишени дизайна долгой жизни поддерживаются опубликованными работами:
	\begin{itemize}
		\item Сверхточная мутация в рибосоме увеличивает продолжительность жизни у дрожжей, червей и мух \cite{Bjedov2021}.
		\item Долгоживущие млекопитающие демонстрируют истощение гипермутабельных кодонов в генах протеостаза \cite{Kerekes2025}.
		\item Сверхэкспрессия митохондриальных шаперонов снижает окислительные повреждения и продлевает жизнь \textit{Drosophila} \cite{Morrow2004}.
		\item Усиленный протеостаз задерживает агрегационно-опосредованную протеотоксичность \cite{Walker2003}.
		\item Возрастная дисрегуляция сигнализации интегринов способствует клеточному старению \cite{Takeda2019}.
	\end{itemize}
	Ни одно из этих исследований, однако, не ввело единый количественный параметр координации. BCLM делает именно это, и многовидовой протокол, приведённый ниже, представляет собой решающую проверку.
	
	\section{Целевые виды и клеточные линии}
	\label{sec:species}
	Были отобраны пять видов, охватывающих широкий филогенетический диапазон при сохранении практических преимуществ клеточного культивирования. Для каждого из них первичные фибробласты являются наиболее доступным нормальным типом клеток, который можно пассировать в течение многих поколений без трансформации.
	
	\begin{table}[H]
		\centering
		\caption{Пять видов и их первичные клеточные линии.}
		\label{tab:species}
		\begin{tabular}{@{}lllll@{}}
			\toprule
			\textbf{Вид} & \textbf{Общее название} & \textbf{Тип клеток} & \textbf{Источник} \\
			\midrule
			\textit{Mus musculus} (C57BL/6) & Домовая мышь & Эмбриональные фибробласты (MEF) & ATCC CRL‑2214 \\
			\textit{Rattus norvegicus} (F344) & Серая крыса & Дермальные фибробласты & Cell Biolabs RAT‑DF \\
			\textit{Sus scrofa} (Landrace) & Домашняя свинья & Дермальные фибробласты & Cell Applications 12405 \\
			\textit{Macaca fascicularis} & Яванский макак & Дермальные фибробласты & Coriell AG08333 \\
			\textit{Canis familiaris} (Beagle) & Собака & Дермальные фибробласты & Coriell AG07090 \\
			\bottomrule
		\end{tabular}
	\end{table}
	
	\section{Ортологи целевых генов и их \(\REL\)}
	\label{sec:genes}
	Для каждого вида мы идентифицировали ортологи человеческих белков, которые являются мишенями в дизайне долгой жизни (BCLM‑LL). Аминокислотные последовательности высококонсервативны; следовательно, вклад аминокислот в \(K_{\mathrm{eff}}\) по существу одинаков у всех видов. Основное различие обусловлено оптимизацией кодонов, поскольку набор оптимальных кодонов зависит от видовой таблицы использования кодонов (база данных Kazusa \cite{codon_usage_db}). Для каждого гена мы вычислили две относительные эффективности координации: \(\REL_{\mathrm{WT}}\) (нативная кодирующая последовательность) и \(\REL_{\mathrm{LL}}\) (кодон-оптимизированная последовательность). Расчёт следовал правилам, описанным в BCLM, раздел~2.
	
	Таблица~\ref{tab:targets} перечисляет ортологи и предсказанные значения \(\REL\) для человеческого референса (из BCLM‑LL) и пяти видов животных. Все значения были сгенерированы скриптом `bclm\_multi\_species.py`, доступным в репозитории YPSDC.
	
	\begin{table}[H]
		\centering
		\caption{Ортологи целевых генов и их предсказанные \(\REL\).}
		\label{tab:targets}
		{\footnotesize
			\begin{tabular}{@{}llcccccc@{}}
				\toprule
				\textbf{Слой} & \textbf{Ген человека} & \textbf{Мышь} & \textbf{Крыса} & \textbf{Свинья} & \textbf{Макак} & \textbf{Собака} & \textbf{Фактор ошибки\footnotemark[1]} \\
				\midrule
				1 & BRCA1 & 0.62→0.78 & 0.63→0.79 & 0.62→0.78 & 0.62→0.78 & 0.85 \\
				& PARP1 & 0.60→0.77 & 0.61→0.78 & 0.60→0.77 & 0.60→0.77 & 0.86 \\
				& MLH1  & 0.59→0.76 & 0.60→0.77 & 0.59→0.76 & 0.59→0.76 & 0.87 \\
				2 & NDUFS1 & 0.55→0.72 & 0.56→0.73 & 0.55→0.72 & 0.55→0.72 & 0.80 \\
				& NDUFV1 & 0.56→0.73 & 0.57→0.74 & 0.56→0.73 & 0.56→0.73 & 0.80 \\
				& NDUFA9 & 0.54→0.70 & 0.55→0.71 & 0.54→0.70 & 0.54→0.70 & 0.81 \\
				3 & HSPA1A & 0.52→0.69 & 0.53→0.70 & 0.52→0.69 & 0.52→0.69 & 0.78 \\
				& DNAJB1 & 0.50→0.68 & 0.51→0.69 & 0.50→0.68 & 0.50→0.68 & 0.79 \\
				& HSPA9  & 0.53→0.70 & 0.54→0.71 & 0.53→0.70 & 0.53→0.70 & 0.78 \\
				4 & ITGAV  & 0.58→0.74 & 0.59→0.75 & 0.58→0.74 & 0.58→0.74 & 0.83 \\
				& ITGB3  & 0.57→0.73 & 0.58→0.74 & 0.57→0.73 & 0.57→0.73 & 0.83 \\
				& FN1    & 0.60→0.76 & 0.61→0.77 & 0.60→0.76 & 0.60→0.76 & 0.84 \\
				\bottomrule
		\end{tabular}}
		\footnotetext[1]{Фактор ошибки \(\varepsilon_{\mathrm{LL}}/\varepsilon_{\mathrm{WT}} = (\REL_{\mathrm{LL}}/\REL_{\mathrm{WT}})^{-2/3}\). Из-за очень высокой консервативности последовательностей фактор варьирует между видами менее чем на \(0.5\%\); приведённые значения являются средними.}
	\end{table}
	
	\section{Предсказанные фенотипические изменения}
	\label{sec:predictions}
	Для каждого слоя общее улучшение получается как среднее геометрическое индивидуальных факторов ошибок белков, поскольку целевые компоненты действуют в перекрывающихся путях. Полученные специфичные для слоёв отношения ошибок:
	
	\begin{itemize}
		\item \textbf{Слой 1 – Стабильность генома:} \(\mu_{\mathrm{LL}}/\mu_{\mathrm{WT}} = \sqrt[3]{0.85\times0.86\times0.87} \approx 0.86\).
		\item \textbf{Слой 2 – Митохондриальные АФК:} \(\mathrm{ROS}_{\mathrm{LL}}/\mathrm{ROS}_{\mathrm{WT}} = \sqrt[3]{0.80\times0.80\times0.81} \approx 0.80\).
		\item \textbf{Слой 3 – Агрегация белков:} \(\mathrm{Agg}_{\mathrm{LL}}/\mathrm{Agg}_{\mathrm{WT}} = \sqrt[3]{0.78\times0.79\times0.78} \approx 0.78\).
		\item \textbf{Слой 4 – Ошибка адгезии к ВКМ:} \(\mathrm{Err}_{\mathrm{LL}}/\mathrm{Err}_{\mathrm{WT}} = \sqrt[3]{0.83\times0.83\times0.84} \approx 0.83\). Поскольку скорость миграции обратно пропорциональна ошибкам адгезии, \(v_{\mathrm{LL}}/v_{\mathrm{WT}} \approx 1/0.83 = 1.20\).
	\end{itemize}
	
	В предположении, что четыре слоя отказывают независимо, общая вероятность фатальной ошибки равна произведению четырёх отношений:
	\[
	\frac{\varepsilon_{\mathrm{fatal,LL}}}{\varepsilon_{\mathrm{fatal,WT}}} \approx 0.86 \times 0.80 \times 0.78 \times 0.83 \approx 0.45.
	\]
	При гипотезе, что репликативная продолжительность жизни (измеряемая кумулятивными удвоениями популяции) обратно пропорциональна вероятности фатальной ошибки, ожидаемое продление жизни составляет
	\[
	\frac{\mathrm{PD}_{\mathrm{LL}}}{\mathrm{PD}_{\mathrm{WT}}} \approx \frac{1}{0.45} \approx 2.2.
	\]
	Это верхняя граница; перекрёстные взаимодействия между слоями и остаточные пути повреждений уменьшат реальный выигрыш. Консервативная оценка, принятая для предварительно зарегистрированного предсказания, помещает фактор продления жизни в диапазон \textbf{1.5–2.0} для всех видов.
	
	Поскольку аминокислотные последовательности почти идентичны, эти предсказания одинаковы для всех пяти видов животных и совпадают с предсказанием для человека из BCLM‑LL.
	
	\section{Экспериментальный протокол}
	\label{sec:protocol}
	
	\subsection{Генетические конструкции}
	Для каждого вида готовятся два синтетических варианта каждого целевого гена:
	\begin{itemize}
		\item \textbf{WT‑OE}: нативная кодирующая последовательность, клонированная в доксициклин-индуцируемый лентивирусный вектор (pLV‑TRE‑Tight‑IRES‑EGFP).
		\item \textbf{LL‑OPT}: оптимизированная BCLM последовательность (кодон-оптимизированная для конкретного вида-хозяина) в том же векторе.
	\end{itemize}
	Для BRCA1 создаётся скремблированный контроль (SCR) путём случайной перестановки синонимичных кодонов.
	
	\subsection{Контроли}
	Для каждого вида параллельно проводятся три контроля:
	\begin{enumerate}
		\item Пустой вектор (EV).
		\item WT‑OE (с согласованными уровнями белка).
		\item SCR (контроль специфичности).
	\end{enumerate}
	
	\subsection{График измерений}
	Для каждого вида выполняются те же пять анализов, что и в BCLM‑LL.
	
	\begin{table}[H]
		\centering
		\caption{Экспериментальный график.}
		\label{tab:schedule}
		\begin{tabular}{@{}llll@{}}
			\toprule
			\textbf{День} & \textbf{Анализ} & \textbf{Слой} & \textbf{Показатель} \\
			\midrule
			0–7   & Трансдукция, селекция & --- & EGFP \\
			7–14  & Индукция доксициклином & Все & уровни белков \\
			14–21 & Анализ мутаций HPRT & 1 & колонии 6‑TG\(^{\mathrm{R}}\) \\
			14–21 & MitoSOX, TMRE & 2 & проточная цитометрия \\
			21–28 & Анализ включений HTT‑Q72 & 3 & флуоресцентная микроскопия \\
			21–28 & Скретч-анализ (рана) & 4 & скорость закрытия раны \\
			28–56 & Серийное пассирование & Все & сенесценс \(\beta\)-gal, теломерная qPCR \\
			\bottomrule
		\end{tabular}
	\end{table}
	
	\subsection{Определение репликативной продолжительности жизни}
	Клетки пассируют при 80\% конфлюэнтности, регистрируют кумулятивные удвоения популяции (CPD). Первичной конечной точкой является CPD при остановке роста. Предсказывается, что оптимизированные BCLM клетки достигнут в \(1.5\)–\(2.0\) раз большего CPD, чем контрольные клетки WT‑OE.
	
	\section{Контролируемая координационная регрессия (CCR) на нескольких видах}
	\label{sec:CCR}
	Обратимость стареющего фенотипа проверяется на каждом виде путём временного подавления одного оптимизированного гена на слой.
	
	\begin{table}[H]
		\centering
		\caption{CCR-вмешательства для пяти видов.}
		\label{tab:CCR}
		\begin{tabular}{@{}llll@{}}
			\toprule
			\textbf{Слой} & \textbf{Целевой ген} & \textbf{Метод} & \textbf{Ожидаемый эффект} \\
			\midrule
			1 & BRCA1\_LL & Dox-индуцируемая shRNA & \(\REL \to 0.62\); частота мутаций возрастает в \(1.86\) раза \\
			2 & NDUFS1\_LL & CRISPR-бейс-редактирование (I182F) & \(\REL \to 0.55\); АФК возрастают в \(1.80\) раза \\
			3 & HSPA1A\_LL & Tet‑CRISPRi & \(\REL \to 0.52\); агрегаты возрастают в \(1.78\) раза \\
			4 & ITGB3\_LL & siRNA & \(\REL \to 0.57\); скорость миграции падает в \(1.20\) раза \\
			\bottomrule
		\end{tabular}
	\end{table}
	
	Через 72 часа подавления биомаркеры должны сместиться в сторону фенотипа дикого типа. После отмывки или разбавления siRNA клетки должны восстановить состояние долгой жизни в течение 48–72 часов. Количественное предсказание для каждого биомаркера выводится непосредственно из значений \(\REL\) в Таблице~\ref{tab:targets} и универсального закона ошибок.
	
	\section{Фальсифицируемость}
	\label{sec:falsifiability}
	Все предсказания в Таблице~\ref{tab:predictions} были предварительно зарегистрированы в репозитории YPSDC. Координационная теория старения YUCT будет считаться опровергнутой, если после трёх независимых повторений экспериментальные результаты для четырёх из пяти видов выйдут за пределы указанных диапазонов.
	
	\begin{table}[H]
		\centering
		\caption{Предварительно зарегистрированные предсказания (одинаковы для всех пяти видов).}
		\label{tab:predictions}
		\begin{tabular}{@{}llll@{}}
			\toprule
			\textbf{Предсказание} & \textbf{Анализ} & \textbf{Ожидаемое (LL/WT)} & \textbf{Опровержение, если} \\
			\midrule
			Частота мутаций       & HPRT           & \(0.86 \pm 0.05\) & LL/WT \(> 0.95\) \\
			АФК                 & MitoSOX        & \(0.80 \pm 0.05\) & LL/WT \(> 0.90\) \\
			Агрегаты          & HTT‑Q72 фокусы   & \(0.78 \pm 0.05\) & LL/WT \(> 0.90\) \\
			Скорость миграции     & Скретч-анализ  & \(1.20 \pm 0.05\) & LL/WT \(< 1.05\) \\
			Репликативная продолжительность жизни & Кумулятивные PD  & \(1.5\)–\(2.0\)   & LL/WT \(< 1.3\) \\
			\bottomrule
		\end{tabular}
	\end{table}
	
	\section{Этические соображения}
	\label{sec:ethics}
	Данный протокол использует исключительно коммерчески доступные клеточные линии. Живые животные не требуются. Генетические конструкции доставляются \textit{in vitro} с помощью стандартных трансфекционных реагентов; методы \textit{in vivo} доставки не описываются. Исследователям настоятельно рекомендуется ознакомиться с Приложением~\(\Sigma\) перед распространением этой работы за пределы культивируемых клеток.
	
	\section{Заключение}
	\label{sec:conclusion}
	Протокол BCLM‑LL‑MPS предоставляет строгую, внутренне согласованную основу для проверки координационной теории старения YUCT на пяти эволюционно различных видах млекопитающих. Все предсказания вычислены из одного набора универсальных констант и видовых таблиц использования кодонов; никакие свободные параметры не подгоняются \textit{post hoc}. Успешный результат установит \(K_{\mathrm{eff}}\) как первый межвидовой, количественный показатель биологической координации, а её восстановление — как рациональную стратегию борьбы со старением.
	
	\vspace{0.5cm}
	\noindent\textbf{Данные и код:} Все ортологичные последовательности, таблицы использования кодонов, вычисленные значения \(\REL\) и предварительно зарегистрированные предсказания доступны по адресу \url{https://github.com/Alexey-Yakushev-YUCT/YPSDC}.
	
	\begin{thebibliography}{99}
		\bibitem{BCLM} A.\,V.\,Yakushev, ``Приложение BCLM: Биологический лагранжиан координации,'' in \emph{YUCT}, Zenodo (2026).
		\bibitem{BCLM_LL} A.\,V.\,Yakushev, ``Приложение BCLM‑LL: Экспериментальный протокол для кардиомиоцитов человека,'' in \emph{YUCT}, Zenodo (2026).
		\bibitem{AppendixL} A.\,V.\,Yakushev, ``Приложение L: Фрактальное масштабирование ошибок координации,'' in \emph{YUCT}, Zenodo (2026).
		\bibitem{AppendixP} A.\,V.\,Yakushev, ``Приложение P: Температурная зависимость гравитации,'' in \emph{YUCT}, Zenodo (2026).
		\bibitem{AppendixSigma} A.\,V.\,Yakushev, ``Приложение \(\Sigma\): Этические императивы и управление,'' in \emph{YUCT}, Zenodo (2026).
		\bibitem{Bjedov2021} I.~Bjedov \emph{et al.}, ``A single hyper‑accurate mutation in the ribosome extends lifespan in yeast, worms, and flies,'' \emph{Cell Metabolism}, 33, 1054–1070 (2021).
		\bibitem{Kerekes2025} K.~Kerekes \emph{et al.}, ``Codon optimisation and evolutionary pressure in long‑lived mammals,'' \emph{bioRxiv} (2025).
		\bibitem{Morrow2004} G.~Morrow \emph{et al.}, ``Overexpression of … Hsp22 extends \textit{Drosophila} lifespan,'' \emph{FASEB J.}, 18, 598–599 (2004).
		\bibitem{Walker2003} G.\,A.~Walker, G.\,J.~Lithgow, ``Lifespan extension in \textit{C. elegans} by a molecular chaperone,'' \emph{Aging Cell}, 2, 131–139 (2003).
		\bibitem{Takeda2019} M.~Takeda \emph{et al.}, ``Downregulation of \(\alpha\)5 integrin induces cellular senescence,'' \emph{FEBS Lett.}, 593, 1284–1293 (2019).
		\bibitem{codon_usage_db} Y.\,Nakamura \emph{et al.}, Codon usage tabulated … \url{https://www.kazusa.or.jp/codon/}.
		\bibitem{YPSDC_repo} A.\,V.\,Yakushev, YPSDC репозиторий, \url{https://github.com/Alexey-Yakushev-YUCT/YPSDC}.
	\end{thebibliography}
	
\end{document}